\section{Ablations And Manual Audits}

\paragraph{PopQA audit.}
The PopQA audit compares a frozen pre-fix bundle against the final post-fix bundle. Before the guardrail fix, the dominant non-recurring failure was: \emph{low-confidence case should have forced retrieval}. After the fix, the remaining failures were mostly: \emph{retrieval was still the right choice}. This is important because it shows that the final PopQA result is not hiding a routing pathology.

\paragraph{Freshness audit.}
The frozen freshness slice has only two failures, so the current manual audit is exhaustive rather than sampled. The two remaining categories are:
\begin{itemize}[leftmargin=1.5em]
\item \emph{retrieval evidence insufficient}: a volatile update where retrieval was the right action class but the visible evidence remained stale,
\item \emph{should have adapted}: a volatile update where the controller fell back to \texttt{PARAM\_ONLY} rather than using temporary adaptation.
\end{itemize}

These audits support a consistent interpretation across both experiments. On PopQA, the final errors are mostly hard retrieval/evidence cases rather than routing mistakes. On freshness, the remaining errors are concentrated in volatile updates rather than in stable confirmation or rollback handling.
